\chapter{Common Mistakes in Imaging}

\section*{Patient focus}
Mistakes in medical imaging can happen, but they are usually preventable. This chapter describes the most common errors that doctors make when reading scans and X-rays, and how to avoid them. Understanding these pitfalls can help patients ask the right questions---for example, ``Has this been compared with my previous scan?'' or ``Is the technical quality of this X-ray adequate?'' If something does not seem right, speaking up can make a difference.

\section*{Imaging focus}
This chapter catalogues the most frequent and high-impact errors in diagnostic imaging. Each mistake includes a description of what goes wrong, how to prevent it, and a clinical example. Recognising these cognitive and technical pitfalls is essential for every clinician who interprets images.
\label{ch:ch100_mistakes}

\section{Why Imaging Errors Happen}
Errors in imaging interpretation arise from three main sources:
\begin{enumerate}
    \item \textbf{Technical factors}---poor image quality, incorrect protocol, incomplete study
    \item \textbf{Cognitive factors}---biases, fatigue, distraction, inexperience
    \item \textbf{System factors}---workload, communication failures, inadequate follow-up
\end{enumerate}

Most critical diagnosis errors are cognitive, not technical. The radiologist sees the abnormality but fails to recognize or communicate its significance.

\clearpage

% ========================================================================
\section{Mistake 1: Diagnosing Cardiomegaly on an AP Film}
% ========================================================================

\subsection{What Goes Wrong}
The cardiothoracic ratio (CTR) is measured on a posteroanterior (PA) chest X-ray. On an anteroposterior (AP) film, the heart is magnified because it is farther from the detector, and a normal heart can appear enlarged. Using the standard CTR $>0.5$ threshold on an AP film leads to overdiagnosis of cardiomegaly.

\subsection{How to Avoid It}
\begin{itemize}
    \item Always check the projection before measuring CTR
    \item On AP films, a CTR up to 0.6 may be normal
    \item If the projection is not labeled, look for the scapulae (retracted from lung fields on PA, overlying lung fields on AP) and the clavicle position
    \item When in doubt, request a PA film or use prior PA films for comparison
\end{itemize}

\subsection{Clinical Example}
A 70-year-old woman is assessed in the emergency department. An AP portable chest X-ray is performed. The report states ``cardiomegaly.'' Subsequent PA erect film shows a normal CTR of 0.45. The patient avoids an unnecessary echocardiogram.

\begin{pearlbox}[Projection matters]
Always determine the projection before measuring the cardiothoracic ratio. An AP film can make a normal heart appear 15--20\% larger than it truly is.
\end{pearlbox}

\clearpage

% ========================================================================
\section{Mistake 2: Missing Pneumothorax on a Supine Film}
% ========================================================================

\subsection{What Goes Wrong}
On a supine chest X-ray, air in the pleural space collects anteriorly and medially rather than at the apex. The classic ``visible pleural line'' at the lung apex may not be present. Instead, signs include deep sulcus sign, increased lucency over the hemithorax, and a sharp heart border.

\subsection{How to Avoid It}
\begin{itemize}
    \item Always be suspicious of pneumothorax in supine trauma patients
    \item Look for the deep sulcus sign (abnormally lucent, depressed costophrenic angle)
    \item Compare with prior films
    \item If clinical suspicion is high and X-ray is equivocal, request CT or lateral decubitus film
\end{itemize}

\subsection{Clinical Example}
A 22-year-old motorcyclist presents after a road traffic collision. Supine CXR is reported as normal. The patient deteriorates with increasing respiratory distress. Repeat CXR (or CT) reveals a large left pneumothorax. The initial film showed subtle deep sulcus sign that was missed.

\clearpage

% ========================================================================
\section{Mistake 3: Confusing a Skin Fold for Pneumothorax}
% ========================================================================

\subsection{What Goes Wrong}
Skin folds (from folds in the patient's skin or clothing) can mimic a pleural line on CXR, leading to a false-positive diagnosis of pneumothorax. Key distinguishing features: a skin fold creates a line with lung markings beyond it, whereas a true pneumothorax shows no lung markings peripheral to the pleural line.

\subsection{How to Avoid It}
\begin{itemize}
    \item Look for lung markings beyond the supposed pleural line
    \item Skin folds often cross anatomical boundaries (e.g., crossing ribs)
    \item The line may be less sharp or less continuous than a true pleural edge
    \item Clinical correlation---is the patient asymptomatic with a CXR done for another reason?
    \item If uncertain, a lateral or oblique view may clarify
\end{itemize}

\subsection{Clinical Example}
An elderly patient undergoing pre-operative assessment has a CXR showing a ``left-sided pneumothorax.'' On closer inspection, faint lung markings are visible beyond the line, which crosses the 5th and 6th ribs---characteristic of a skin fold. No pneumothorax is present, and surgery proceeds as planned.

\clearpage

% ========================================================================
\section{Mistake 4: Forgetting to Check Technical Quality}
% ========================================================================

\subsection{What Goes Wrong}
Interpreting an image before assessing its technical adequacy leads to missed findings and erroneous conclusions. A poorly positioned, overexposed, or motion-blurred image may hide pathology or create artefacts that mimic disease.

\subsection{How to Avoid It}
Make technical assessment a habit---every image, every time:
\begin{enumerate}
    \item \textbf{CXR}---Is it PA or AP? Full inspiration? Rotation? Penetration adequate?
    \item \textbf{CT}---IV contrast given? Appropriate phase? Complete coverage of the region of interest?
    \item \textbf{MRI}---Appropriate sequences obtained? Adequate signal-to-noise? Motion artefact?
    \item \textbf{Ultrasound}---Appropriate probe? Full examination of the structure in two planes?
\end{enumerate}

\begin{pearlbox}[A bad image is worse than no image]
A technically limited study with a normal report can create false reassurance. Always state the technical adequacy before interpreting findings. If the study is suboptimal, recommend repeating it.
\end{pearlbox}

\clearpage

% ========================================================================
\section{Mistake 5: Anchoring Bias}
% ========================================================================

\subsection{What Goes Wrong}
Anchoring bias occurs when the first piece of information received (e.g., a referral diagnosis, a prior report) disproportionately influences interpretation. The radiologist ``anchors'' to the initial hypothesis and fails to consider alternative diagnoses.

\subsection{How to Avoid It}
\begin{itemize}
    \item Read the clinical question, but form your own impression before reading the referral diagnosis
    \item Systematically search all structures on the image, not just the area of interest
    \item Ask: ``What else could explain these findings?''
    \item Use structured reporting templates to force comprehensive evaluation
\end{itemize}

\subsection{Clinical Example}
A CT head is requested with the clinical note ``?stroke.'' The radiologist identifies a small hypodensity in the territory of the middle cerebral artery and reports ``acute infarct.'' The posterior fossa, which contains a cerebellar metastasis, is not carefully evaluated because attention was anchored to the supratentorial region.

\clearpage

% ========================================================================
\section{Mistake 6: Satisfaction of Search}
% ========================================================================

\subsection{What Goes Wrong}
After identifying one abnormality, the search for additional findings stops. This is one of the most well-documented cognitive biases in radiology. Studies show that when a radiologist finds one lesion, the detection rate for a second lesion drops by up to 30\%.

\subsection{How to Avoid It}
\begin{itemize}
    \item Adopt a systematic search pattern for every study (e.g., ABCDE for CXR, compartments for CT head)
    \item After identifying an abnormality, consciously restart your search from the beginning
    \item Use checklists or structured reporting to ensure complete evaluation
    \item If you find one lesion, specifically ask: ``Is there another?''
\end{itemize}

\subsection{Clinical Example}
A mammogram shows a suspicious mass in the right breast. The radiologist focuses on this finding and recommends biopsy. A second, smaller satellite lesion in the same breast---which would change staging---is missed. A systematic two-breast, all-quadrant search pattern would have identified both lesions.

\begin{pearlbox}[Never stop after the first finding]
Satisfaction of search is the single most common cause of double pathology being missed. After finding one abnormality, restart your search systematically.
\end{pearlbox}

\clearpage

% ========================================================================
\section{Mistake 7: Missing Free Air Under the Diaphragm}
% ========================================================================

\subsection{What Goes Wrong}
Free intraperitoneal air (pneumoperitoneum) is a surgical emergency indicating perforated viscus. On an upright CXR or abdominal X-ray, it appears as a radiolucent crescent beneath the diaphragm. It is easily missed if the clinician does not specifically look for it, or if only a supine film is obtained.

\subsection{How to Avoid It}
\begin{itemize}
    \item Always examine the area under both hemidiaphragms on every upright CXR
    \item Free air is visible on an upright film with as little as 1--2 mL of intraperitoneal gas
    \item A left lateral decubitus position (right side up) is more sensitive for small amounts of free air
    \item If clinical suspicion is high and plain films are negative, request CT abdomen (most sensitive)
\end{itemize}

\subsection{Clinical Example}
A 68-year-old man with a history of peptic ulcer disease presents with sudden-onset epigastric pain. An upright CXR shows a thin crescent of lucency under the right hemidiaphragm. This is free air from a perforated duodenal ulcer, and the patient proceeds to emergency surgery.

\clearpage

% ========================================================================
\section{Mistake 8: Confusing Appendicitis with a Normal Appendix}
% ========================================================================

\subsection{What Goes Wrong}
On CT, a normal appendix can be seen in up to 70\% of studies. A normal appendix is $<6$ mm in diameter, has a thin wall, and has no surrounding fat stranding. Overcalling a normal or mildly prominent appendix as appendicitis leads to unnecessary surgery.

\subsection{How to Avoid It}
\begin{itemize}
    \item Apply strict CT criteria: appendix $>6$ mm WITH wall thickening AND periappendiceal fat stranding
    \item A dilated appendix without secondary signs may be a normal variant or represent early/resolving appendicitis
    \item If the appendix is visible, non-dilated, and thin-walled, it is normal---regardless of clinical suspicion
    \item Clinical correlation is essential: the Alvarado score should guide the decision to operate
\end{itemize}

\begin{pearlbox}[Visible normal appendix = appendicitis excluded]
A non-dilated, thin-walled appendix that is fully visualised on CT effectively excludes appendicitis, even in a patient with right iliac fossa pain. Do not operate on imaging findings alone without clinical correlation.
\end{pearlbox}

\clearpage

% ========================================================================
\section{Mistake 9: Overcalling Incidental Findings}
% ========================================================================

\subsection{What Goes Wrong}
Incidental findings---abnormalities discovered while investigating an unrelated condition---are extremely common. Overcalling clinically insignificant findings leads to unnecessary investigations, patient anxiety, and healthcare costs. Examples include simple renal cysts, hepatic haemangiomas, benign bone islands, and thyroid nodules.

\subsection{How to Avoid It}
\begin{itemize}
    \item Learn the imaging features of common benign incidental findings
    \item Follow established guidelines (Fleischner for lung nodules, TI-RADS for thyroid, Bosniak for renal cysts)
    \item Do not investigate findings that meet criteria for benignity
    \item Communicate incidental findings clearly, distinguishing them from the primary finding
\end{itemize}

\subsection{Clinical Example}
A CT abdomen for suspected appendicitis identifies a 2 cm simple renal cyst (Bosniak I). The report appropriately notes this as a benign incidental finding with no follow-up required, avoiding unnecessary renal ultrasound and patient concern.

\clearpage

% ========================================================================
\section{Mistake 10: Not Comparing with Prior Studies}
% ========================================================================

\subsection{What Goes Wrong}
Interpreting an imaging study in isolation, without comparison to prior examinations, dramatically reduces diagnostic accuracy. Change over time is one of the most important clues to the nature of a finding---a stable nodule over 2 years is almost certainly benign, while a rapidly growing mass is highly suspicious.

\subsection{How to Avoid It}
\begin{itemize}
    \item Always request prior studies for comparison before issuing a report
    \item Document whether prior studies were available and what comparison was made
    \item When prior studies are unavailable, state this explicitly in the report
    \item Use PACS or enterprise imaging systems to automatically retrieve prior studies
\end{itemize}

\subsection{Clinical Example}
A chest X-ray shows a 1.5 cm pulmonary nodule. Without prior imaging, this requires follow-up per Fleischner criteria (CT at 6--12 months). However, comparison with a CXR from 3 years ago reveals the nodule was the same size then---stability over 2 years in a solid nodule $<8$ mm essentially excludes malignancy, and no follow-up is needed.

\begin{pearlbox}[Stability is the best test]
For many imaging findings (pulmonary nodules, bone lesions, thyroid nodules), stability over time is the strongest evidence of benignity. Always compare with prior studies.
\end{pearlbox}

\clearpage

% ========================================================================
\section{Mistake 11: Missing a Fracture on Initial Imaging}
% ========================================================================

\subsection{What Goes Wrong}
Subtle fractures---particularly occult hip fractures, scaphoid fractures, and rib fractures---may be invisible on initial X-rays. Failing to follow up clinically suspicious cases with advanced imaging (MRI or CT) leads to missed diagnoses and complications.

\subsection{How to Avoid It}
\begin{itemize}
    \item If clinical suspicion remains high despite a normal X-ray, proceed to MRI (most sensitive) or CT
    \item For suspected scaphoid fracture with normal X-ray---immobilise and repeat X-ray at 10--14 days, or MRI immediately
    \item For elderly patients with hip pain and normal X-ray---MRI within 48 hours
    \item Document the clinical reason for advanced imaging when initial X-ray is negative
\end{itemize}

\clearpage

% ========================================================================
\section{Mistake 12: Inadequate Communication of Critical Findings}
% ========================================================================

\subsection{What Goes Wrong}
A critical finding (e.g., tension pneumothorax, aortic dissection, intracranial hemorrhage) is identified on imaging but not communicated to the clinical team in a timely manner. This can occur due to phone unavailability, incorrect contact details, assumption that the clinician will see the report, or failure to escalate.

\subsection{How to Avoid It}
\begin{itemize}
    \item Know your institution's critical results policy
    \item Document the name of the person notified, time, and content of communication
    \item If unable to reach the clinical team, escalate to the department lead
    \item Use structured critical result communication (SBAR format)
    \item Do not rely solely on the electronic report for critical findings
\end{itemize}

\begin{pearlbox}[A finding is not communicated until you confirm the other person heard it]
Telephone or verbal notification of critical findings must be documented with the recipient's name. An unsigned report sitting in a worklist is not communication.
\end{pearlbox}

\clearpage

% ========================================================================
\section{Mistake 13: Forgetting to Check for Prior Contrast Reactions}
% ========================================================================

\subsection{What Goes Wrong}
Contrast safety screening should identify prior reactions to the same contrast class, relevant kidney risk, and patient-specific factors. Severe reactions are uncommon, and kidney screening practices depend on the agent, route, clinical context and current guidance.

\subsection{How to Avoid It}
\begin{itemize}
    \item Ask about prior reactions to the same contrast class and document the agent and reaction when known; unrelated allergies alone do not establish a contrast allergy.
    \item Follow current, agent-specific renal screening guidance. Consider acute kidney injury and severe chronic kidney disease; do not assume every patient needs an eGFR test before every contrast study.
    \item After a prior allergic-like reaction, consider an alternative agent or modality and follow institutional guidance. Premedication does not eliminate risk and is not a substitute for a risk-benefit decision.
    \item For gadolinium, use current ACR risk groups. A clinically indicated Group II agent is generally not withheld solely for eGFR $<30$; assess the individual situation and local policy.
\end{itemize}

\clearpage

% ========================================================================
\section{Mistake 14: Misidentifying Laterality}
% ========================================================================

\subsection{What Goes Wrong}
Reporting pathology on the wrong side---a left-sided finding described as right, or vice versa---is a never event. It can occur due to mirror-image display, incorrect labelling, or simple human error.

\subsection{How to Avoid It}
\begin{itemize}
    \item Always confirm laterality before reporting
    \item Check the image orientation markers (R/L markers on X-ray, digital orientation tags on CT/MRI)
    \item Use standardised reporting templates that explicitly state laterality
    \item If laterality is uncertain, state this and recommend clarification
\end{itemize}

\clearpage

% ========================================================================
\section{Mistake 15: Failure to Recommend Follow-Up}
% ========================================================================

\subsection{What Goes Wrong}
Identifying a finding that requires interval follow-up (e.g., a pulmonary nodule, a renal mass, a bone lesion) but failing to document a follow-up recommendation or ensure that the patient is aware.

\subsection{How to Avoid It}
\begin{itemize}
    \item Every report should end with a clear recommendation for follow-up if applicable
    \item Use structured templates that include a ``recommendation'' field
    \item For critical follow-up (e.g., suspicious nodule), communicate directly with the referring clinician
    \item Consider automated follow-up tracking systems in PACS/RIS
\end{itemize}

\begin{pearlbox}[A finding without a plan is a liability]
If you identify a finding that needs follow-up, document: what to do, when to do it, and who is responsible. A report that says ``recommend follow-up'' without specifying the modality, interval, or responsible clinician is incomplete.
\end{pearlbox}

\clearpage

\section{Summary}
The most dangerous errors in imaging are cognitive, not technical. Anchoring bias, satisfaction of search, and failure to compare with prior studies account for a large proportion of missed diagnoses. Prevention requires systematic search patterns, structured reporting, clinical correlation, and a culture where checking and double-checking is valued over speed. Every radiologist---from trainee to consultant---is susceptible to these biases. Awareness is the first step; habits and systems are the defence.

\clearpage
